\scanpage{178}
fortiori $\underline O_{X^{p/S}}\xrightarrow{\sim}\Lambda$ \deleted{\unsure{Donc On sait bi|Donc On sait bien}} [Il suffit de vérifier fibre par fibre\ldots\ \inserted{or $X_s\longrightarrow X_s$ est fini, $X_s$ (\unsure{abst|étant}) régulier, \unsure{$=$ dim pour source|$=$ dim que source|à dimension plus source}.}] On sait bien que $X/S$ a localmt une $p$-base i.e. qu'il en est ainsi de $X$ sur $X^{p/S}$ : \unsure{en effet|enf} fait, grâce au fait que $X$ est loc de t.f. sur $S$, cette condition équivaut à celle que $\Omega^1_{X/S}$ soit localement libre, ce qui est bien le cas !

On est donc dans un cas où l'opération de Cartier est
\[
 \widetilde C:f_{X/S*}^{\deleted{\unsure{*}}}(Z\Omega^*_{X/S})\longrightarrow\Omega^*_{X^{p/S}}
\]
Si la dim relative est $n$, on trouve en particulier
\[
 \widetilde C:f_{X/S*}(\underbrace{\underline\Omega^n_{X/S}}_{
 \substack{\Vert\\\omega_{X/S}\\\Vert\\\text{module}\\\text{dualisant}\\\text{relatif}}})
 \longrightarrow\underset{\substack{\Vert\\\omega_{X^{p/S}}}}{\underline\Omega^n_{X^{p/S}}}
\]

\uline{Théorème} 3\quad L'homomorphisme précédent est l'homomorphisme trace (ou de Gysin) $f_{X/S*}$ associé au morphisme $f_{X/S}:X\longrightarrow X^{p/S}$ de schémas lisses sur $S$.

Soit alors $f_1,\ldots,f_n$ une suite de sections de $\underline O_X$ qui définissent un sous-schéma fermé \inserted{$Z$} de $X$ \uline{fini} sur $S$ \inserted{(donc loc lib. sur $S$)}. $\widetilde f_1\ldots\widetilde f_n$ définissent donc $\widetilde Z$ itou, et on peut former, pour une section $\overline\omega$ de $\Omega^n_{X^{p/S}}$, $\deleted{\operatorname{Tr}}\operatorname{res}\left(\frac{\overline\omega}{\widetilde f_1\ldots\widetilde f_n}\right)$. On peut aussi regarder les images $f_1^p,\ldots,f_n^p$ dans $X$, qui définissent un sous-schéma $Z'$ de $X$ fini sur $S$, et pour les $n$-formes $\omega$ de $X/S$, on peut former $\operatorname{res}\left(\frac{\omega}{f_1^p\ldots f_n^p}\right)$. On trouve alors, comme façon essentielle
% archive sheet number: 177

